\chapter{Styles, Style Options, and Delegates}

Qt styles are the part of Qt that decide how standard controls look and how much space they need. A push button, combo box, progress bar, tab, toolbar button, focus rect, or item-view row is not painted by hand every time you use one. Qt asks the active style to do that work.

That matters for two reasons:

\begin{enumerate}
  \item Most applications should let the current style do as much work as possible so controls continue to look native or at least coherent.
  \item When you do need to customize presentation, the style APIs give you a lower-level way to stay inside Qt's layout and paint rules instead of replacing everything with manual drawing.
\end{enumerate}

The style surface in \api{crystal-qt6} now covers three different levels:

\begin{longtable}{@{}p{0.28\textwidth}p{0.66\textwidth}@{}}
\toprule
Type & Use it for \\
\midrule
\api{Style}, \api{CommonStyle}, \api{StyleFactory} & Selecting the application or widget look-and-feel and working with standard palettes. \\
\api{StyleOption*} and \api{StylePainter} & Low-level, style-aware drawing when you need Qt to render a control, title bar, tab, progress bar, or item-view row into a custom target. \\
\api{StyledItemDelegate} and item editor factories & Customizing how model/view data is displayed, sized, painted, and edited without abandoning the view's style behavior. \\
\bottomrule
\end{longtable}

\section{Choosing The Right Tool}

There are several ways to influence a Qt interface, and they solve different problems.

\begin{longtable}{@{}p{0.24\textwidth}p{0.24\textwidth}p{0.42\textwidth}@{}}
\toprule
Tool & Best for & Notes \\
\midrule
\api{StyleFactory} or \api{Application#style=} & Picking an overall application look, such as the current platform style or Fusion. & This affects standard widget drawing, spacing, metrics, and many default colors. It is the first lever to reach for. \\
\api{QPalette} & Adjusting theme-aware colors without replacing the control chrome. & Good for custom dark/light palettes, warning accents, or project color systems that should still feel native. \\
\api{style\_sheet} & Targeted visual overrides for a small number of widgets. & Useful, but easy to overuse. Large style sheets can pull the UI away from the active style and make behavior less predictable across platforms. \\
\api{StylePainter} and \api{StyleOption*} & Rendering a style-controlled control into an image, export target, or custom widget paint path. & Use this when you want Qt to draw a control for you, but not as a live widget. \\
\api{StyledItemDelegate} & Custom list, tree, and table presentation and editing. & This keeps item views style-aware while still letting you format text, size rows, paint cells, or choose editors. \\
\bottomrule
\end{longtable}

If the question is ``how do I make the whole application feel different?'' use a style or palette first. If the question is ``how do I draw one control-shaped thing correctly?'' use a style option plus \api{StylePainter}. If the question is ``how do I customize rows in a view?'' use \api{StyledItemDelegate}.

\begin{figure}[htbp]
  \centering
  \screenshotimage{docs/book/images/styles-workbench.png}{width=0.92\textwidth,height=0.58\textheight,keepaspectratio}%
  \caption{A small styles workbench: standard controls under the Fusion style, a styled painter preview, and a delegate-backed editor row.}
\end{figure}

\section{Selecting A Style}

The simplest style change is to ask Qt for the styles available on the current platform, then pick one. Fusion is the most predictable cross-platform choice when you want screenshots, tutorials, or tests to look reasonably consistent.

\begin{lstlisting}[style=crystal]
app = Qt6.application

if Qt6::StyleFactory.keys.includes?("Fusion")
  app.set_style("Fusion")
end
\end{lstlisting}

You can also create a style instance directly and assign it:

\begin{lstlisting}[style=crystal]
app = Qt6.application
fusion = Qt6::StyleFactory.create("Fusion")

unless fusion.nil?
  app.style = fusion
end
\end{lstlisting}

Assigning a style to the application is the normal path. Assigning one to a specific widget is possible, but should usually be the exception:

\begin{lstlisting}[style=crystal]
inspector = Qt6::Widget.new
local_style = Qt6::StyleFactory.create("Fusion")

unless local_style.nil?
  inspector.style = local_style
end
\end{lstlisting}

Widget-local styles are most useful when one panel needs an intentionally different presentation, or when you are experimenting. For production applications, a single application style plus palette adjustments is usually easier to keep coherent.

\section{Styles Versus Style Sheets}

Qt style sheets are still useful, but they solve a different problem than \api{QStyle}. A style picks the platform-aware drawing logic and metrics. A style sheet overrides widget appearance using a CSS-like layer.

That means:

\begin{itemize}
  \item A style decides how a button, tab, or progress bar is normally drawn.
  \item A palette adjusts the colors the style uses.
  \item A style sheet can override those decisions for a specific widget or subtree.
\end{itemize}

For broad desktop polish, prefer style plus palette. Reach for style sheets when you want a clearly intentional local treatment, such as a warning banner, a canvas-side inspector, or a branded splash screen. The more of the application you restyle with CSS, the less helpful the active Qt style becomes.

\section{Style-Aware Painting}

Sometimes you want a control to be drawn by Qt, but not as a live widget in a layout. This is what \api{StylePainter} and \api{StyleOption*} are for.

A style option is a structured packet of state: geometry, enabled or selected flags, text, icon, orientation, focus state, and so on. A style painter uses that packet to ask the current style to draw the control.

\begin{lstlisting}[style=crystal]
app = Qt6.application
button = Qt6::PushButton.new("Deploy")
button.resize(140, 36)
button.show
app.process_events

option = Qt6::StyleOptionButton.new
option.init_from(button)
option.rect = Qt6::Rect.new(0, 0, 140, 36)

image = Qt6::QImage.new(180, 52)
image.fill(Qt6::Color.new(250, 250, 250))

painter = Qt6::StylePainter.new
if painter.begin(image, button)
  painter.draw_control(Qt6::StyleControlElement::PushButton, option)
  painter.end
end

image.save("deploy-button.png")
\end{lstlisting}

This is the same pattern you use for other controls:

\begin{longtable}{@{}p{0.28\textwidth}p{0.66\textwidth}@{}}
\toprule
Option type & Common use \\
\midrule
\api{StyleOptionButton} & Paint a standard push button or inspect button state. \\
\api{StyleOptionProgressBar} & Render a progress bar into an offscreen image or export target. \\
\api{StyleOptionSlider} & Paint slider-like controls or inspect their ranges and handle state. \\
\api{StyleOptionTab} and \api{StyleOptionTabBarBase} & Style-aware tab drawing and tab geometry. \\
\api{StyleOptionToolButton} and \api{StyleOptionToolBar} & Tool palettes, toolbars, and menu-button chrome. \\
\api{StyleOptionViewItem} & Item-view rows, cells, and delegate paint state. \\
\bottomrule
\end{longtable}

Most applications do not create every option family directly. Treat them as low-level interoperability types: use the narrowest one that matches the control you need to paint.

\begin{figure}[htbp]
  \centering
  \screenshotimage{docs/book/images/styles-style-painter.png}{width=0.82\textwidth,height=0.30\textheight,keepaspectratio}%
  \caption{One offscreen image rendered with \api{StylePainter}: a button, combo box, progress bar, focus rect, text, and a status badge, all drawn by the active style.}
\end{figure}

\section{Delegates And Editors}

In model/view code, styles matter most through delegates. \api{StyledItemDelegate} gives you a style-aware way to customize display text, item paint behavior, size hints, and editor creation without replacing the view itself.

\begin{lstlisting}[style=crystal]
delegate = Qt6::StyledItemDelegate.new(table)

delegate.on_display_text do |text|
  "#{text}  [reviewed]"
end

delegate.on_create_editor do |parent, _index|
  Qt6::LineEdit.new(parent: parent)
end

delegate.on_set_editor_data do |editor, value, _index|
  editor.as(Qt6::LineEdit).text = value.to_s
end

delegate.on_set_model_data do |editor, model, index|
  model.set_data(index, editor.as(Qt6::LineEdit).text)
end

table.item_delegate = delegate
\end{lstlisting}

This is the main delegate progression:

\begin{enumerate}
  \item Use model roles when you only need different text, icons, check state, colors, or tooltips.
  \item Add \api{StyledItemDelegate#on\_display\_text} when the presentation is mostly text formatting.
  \item Add editor hooks when the default editor choice is not enough.
  \item Reach for \api{on\_paint}, \api{on\_size\_hint}, or direct \api{paint} and \api{size\_hint} helpers when row layout or painting itself needs to change.
\end{enumerate}

The wrapper now also exposes the delegate's direct public helpers, which are useful for tests, previews, and export flows:

\begin{lstlisting}[style=crystal]
option = delegate.init_style_option(index)
size = delegate.size_hint(option, index)

Qt6::QPainter.paint(image) do |painter|
  delegate.paint(painter, option, index)
end
\end{lstlisting}

\section{Editor Factories}

Qt can also centralize editor selection through \api{QItemEditorFactory}. This is most useful when several delegates or several views should create the same editor type for the same kind of value.

\begin{lstlisting}[style=crystal]
factory = Qt6::QItemEditorFactory.new
creator = Qt6::QStandardItemEditorCreator.new

creator.on_create_widget do |parent|
  Qt6::LineEdit.new(parent: parent)
end

factory.register_editor_for("terrain", creator)
delegate.set_item_editor_factory(factory)
\end{lstlisting}

In practice, the item editor factory is an optional refinement. Many applications do perfectly well with explicit delegate hooks and never need a shared factory layer.

\section{Proxy Styles And Plugin Hooks}

\api{ProxyStyle} and \api{StylePlugin} are advanced tools, and it is worth setting expectations clearly.

\api{ProxyStyle} currently gives you a way to wrap or replace a base style object:

\begin{lstlisting}[style=crystal]
base_style = Qt6::StyleFactory.create("Fusion")
proxy = Qt6::ProxyStyle.new(base_style)
\end{lstlisting}

That is useful for lifetime management and for keeping a style object attached to a widget or application boundary, but the Crystal wrapper does not yet expose Qt's full virtual override surface for custom proxy-style behavior. Think of it today as style composition infrastructure, not yet as a full style-subclassing API.

\api{StylePlugin} is similarly an advanced hook. It is useful when you want a callback-backed style factory inside Crystal, but it is not the first tool to reach for when the real goal is just ``make the app look different.'' Most readers can safely ignore it until they are already comfortable with regular styles, palettes, and delegates.

\section{Practical Guidance}

When in doubt:

\begin{itemize}
  \item Start with the active platform style or Fusion.
  \item Use palettes before broad style sheets.
  \item Use style-aware drawing when exporting or previewing controls offscreen.
  \item Use delegates for view customization instead of painting directly into the viewport.
  \item Treat the large \api{StyleOption*} family as a low-level toolbox, not as a surface you must learn all at once.
\end{itemize}

That approach keeps your application easier to reason about and lets Qt keep doing the hard work of spacing, metrics, state handling, and platform-appropriate control drawing.
